% Irodov Part III §3.4 Electric Current, 147--160;
% PDF pp.121--123, printed pp.112--114; figure 3.35--3.38 checked.
\electrocurrentproblem{147}{112}
半径为 $a=1.0\,\mathrm{cm}$ 的无限长圆柱
侧面均匀带电，并沿自身轴线以
$v=10\,\mathrm{m/s}$ 匀速平动。
圆柱表面外电场强度为
$E=0.9\,\mathrm{kV/cm}$。
求电荷随物质平动产生的传导外形电流（对流电流）。
\exsolution
无限均匀线电荷的表面外场
$E=\lambda/(2\pi\varepsilon_0a)$，
故单位轴长电荷量
$\lambda=2\pi\varepsilon_0aE$。
圆柱每秒沿轴移动长度 $v$，
每秒通过固定截面的电荷
$I=\lambda v$，所以
\[
\boxed{I=2\pi\varepsilon_0aEv}.
\]
代入 $a=0.010\,\mathrm m$、
$E=0.9\,\mathrm{kV/cm}
=9.0\times10^4\,\mathrm{V/m}$：
\[
\boxed{I\simeq5.0\times10^{-7}\,\mathrm A
=0.50\,\mathrm{\mu A}}.
\]
\excheck 这是带电物质整体移动产生的\emph{对流电流}，
并非静止金属内自由电子的漂移电流。

\electrocurrentproblem{148}{112}
电压固定为 $V=200\,\mathrm V$ 的空气
同轴圆柱电容器以速度 $v=5.0\,\mathrm{m/s}$
竖直浸入水中，径向间隙
$d=2.0\,\mathrm{mm}$、平均半径
$r=50\,\mathrm{mm}$，$d\ll r$。
求电源导线中的电流模。
\exsolution
浸水长度增加 $\mathrm dx=v\,\mathrm dt$。
沿圆柱轴每单位长度的真空电容为
$C'_0\simeq2\pi\varepsilon_0r/d$。
被水替换时增量每单位长度为
$(\varepsilon-1)C'_0$，
因此电容增长率
\[
\frac{\mathrm dC}{\mathrm dt}
=\frac{2\pi\varepsilon_0(\varepsilon-1)rv}{d}.
\]
电压固定，电荷 $Q=CV$，
源电流
\[
\boxed{I=V\frac{\mathrm dC}{\mathrm dt}
=\frac{2\pi\varepsilon_0(\varepsilon-1)rvV}{d}}.
\]
室温水取 $\varepsilon\simeq81$，
代入原印速度 $v=5.0\,\mathrm{m/s}$，
\[
\begin{aligned}
\frac{\mathrm dC}{\mathrm dt}
&=\frac{2\pi(8.854\times10^{-12})(80)
(0.050)(5.0)}{0.0020}
\\
&\simeq5.56\times10^{-7}\,\mathrm{F/s},
\\
I&=(200\,\mathrm V)(5.56\times10^{-7}\,\mathrm{F/s})
\\
&\simeq1.11\times10^{-4}\,\mathrm A.
\end{aligned}
\]
故
\[
\boxed{I\simeq0.11\,\mathrm{mA}
=110\,\mathrm{\mu A}}.
\]
\excheck 原版英文扫描本题确实印速度
$v=5.0\,\mathrm{m/s}$；其书末简答却写
$0.11\,\mathrm{\mu A}$，与原公式、单位换算相差
$10^3$ 倍。若速度原拟为 $5.0\,\mathrm{mm/s}$
才会得到该书末数值，但\emph{不能擅自把原印
$\mathrm{m/s}$ 改写成 $\mathrm{mm/s}$}。
所以本译稿保留题给速度，修正结果并在勘误记录说明。
电流来自电容随浸水长度变化，不是水的欧姆漏电流。

\electrocurrentproblem{149}{112--113}
在 $0^\circ\mathrm C$ 时，两导体电阻
$R_{2,0}=\eta R_{1,0}$，
温度系数为 $\alpha_1,\alpha_2$。
求串联与并联连接后电阻的一阶温度系数。
\exsolution
按低温差的一阶线性模型
$R_i(t)=R_{i,0}(1+\alpha_i t)$。
(a) 串联：
$R_s(t)=R_{1,0}+R_{2,0}
+t(R_{1,0}\alpha_1+R_{2,0}\alpha_2)$，
所以
\[
\boxed{\alpha_s=
\frac{\alpha_1+\eta\alpha_2}{1+\eta}}.
\]
(b) 并联电导相加，
$G_p=1/R_1+1/R_2$，
其导数在 $t=0$ 为
$G_p'(0)=-(\alpha_1/R_{1,0}
+\alpha_2/R_{2,0})$，
并有
$\alpha_p=-G_p'/G_p$，因此
\[
\boxed{\alpha_p=
\frac{\eta\alpha_1+\alpha_2}{1+\eta}}.
\]
\excheck 只有电阻温度变化\emph{一阶}时才有
上述简单权重；串联权重是电阻，
并联权重是电导。

\electrocurrentproblem{150}{113}
立方体金属线框，每条棱电阻均为 $R$，
顶点编号见图3.35。
求两点间等效电阻：(a) 1—7；
(b) 1—2；(c) 1—3。
\exsolution
将每条棱看作 $R$ 的电阻，
以相关顶点加电压 $V$，
其它六个顶点应用 Kirchhoff 节点定律。
(a) 1、7 是体对角相对的两个顶点。
由立方体三重对称性，离1一条棱的3顶点同势，
离1两条棱的另3顶点同势。
网络化为
$R/3+R/6+R/3$ 串联，
因此
\[
\boxed{R_{1,7}=\frac56R}.
\]
(b) 1、2 相邻。把两个端点电势分别设为
$V,0$，按垂直二等分面对称对其它节点
写节点方程并消元，可得源注入电流
$I=12V/(7R)$，
所以
\[
\boxed{R_{1,2}=\frac7{12}R}.
\]
(c) 1、3 在同一面的对角位置。
关于连接两点的对称变换，将两个
与端点等距的顶点合并，再使用节点方程，
所得源电流 $I=4V/(3R)$，故
\[
\boxed{R_{1,3}=\frac34R}.
\]
\excheck 结果满足
$\frac7{12}R<\frac34R<\frac56R$，
与相邻、面对角、体对角的几何距离排序一致。

\electrocurrentproblem{151}{113}
图3.36为有限节重复梯形电阻网络，
水平各段为 $2R$，
各节向下支路为 $R$，
右端连接终端电阻 $R_x$。
求让 A、B 端等效电阻与重复节数无关的 $R_x$。
\exsolution
从一个纵向分支顶端向下观察，
无限重复网络的等效电阻记为 $X$。
本节点电阻 $R$ 与其右边
“水平 $2R$ 串接另一份相同网络 $X$”
并联：
\[
X=R\parallel(2R+X)
=\frac{R(2R+X)}{3R+X}.
\]
整理：
$X(3R+X)=R(2R+X)$，
即
\[
X^2+2RX-2R^2=0.
\]
取正根：
\[
\boxed{R_x=X=(\sqrt3-1)R}.
\]
\excheck 将最右端网络用此固定点等效电阻替代，
每向左新增一节仍维持同一个 $X$，
故 A、B 端等效值不随总节数改变。

\electrocurrentproblem{152}{113}
图3.37的无穷梯形网络，
每节水平电阻
$R_1=4.0\,\Omega$，
每节向下电阻
$R_2=3.0\,\Omega$。
求从 A 到 B 的等效电阻。
\exsolution
设整个网络等效为 $X$。
首段为 $R_1$ 串联
$R_2$ 与余下相同网络 $X$ 的并联：
\[
X=R_1+\frac{R_2X}{R_2+X}.
\]
消去分母：
$X(R_2+X)=R_1(R_2+X)+R_2X$，
于是
$X^2-R_1X-R_1R_2=0$。
正根为
\[
\boxed{X=\frac{R_1}{2}
\left(1+\sqrt{1+\frac{4R_2}{R_1}}\right)}.
\]
代入 $R_1=4,R_2=3$：
\[
\boxed{X=2(1+\sqrt4)=6\,\Omega}.
\]
\excheck 负根对应非物理的被动电阻；
递推解与每段串并联拓扑一致。

\electrocurrentproblem{153}{113}
图3.38为无限大平面正方形电阻网格，
任意两相邻交点之间导线的电阻均为
$R_0$。求相邻节点 A、B 间的等效电阻。
\exsolution
采用对称性与电流叠加。
先令电流 $I$ 从 A 注入，
在无限远均匀流出。由于网格在 A 点
的四个方向等价，
从 A 直接流入 A—B 导线的电流为
$I/4$。
再令同样电流从无限远均匀流入，
在 B 点流出，由 B 点四向对称，
A—B 导线从 A 到 B 的电流也是 $I/4$。
线性叠加这两种分布，
得到 A 注入、B 流出的实际电流，
其中 A—B 这条直接边的电流
$I_{AB}=I/2$。
由于 A、B 间电势差
$V_{AB}=I_{AB}R_0$，
\[
\boxed{R_{\rm eff}=\frac{V_{AB}}{I}
=\frac{R_0}{2}}.
\]
\excheck 关键是使用\emph{无限大}
方格网和 A、B 相邻的四重对称；
有限网格的边界会破坏精确的 $R_0/2$。

\electrocurrentproblem{154}{113--114}
无限接近同轴的两根理想金属圆柱导体，
半径 $a<b$、长度 $\ell$，
其间填充均匀弱导电介质，电阻率 $\rho$。
忽略端部效应，求两筒间电阻。
\exsolution
电流沿半径方向，
半径为 $r$ 的圆柱面导电面积为
$A(r)=2\pi r\ell$，
厚度 $\mathrm dr$ 的薄层电阻
$\mathrm dR=\rho\,\mathrm dr/(2\pi r\ell)$。
各径向薄层串联，积分：
\[
R=\frac{\rho}{2\pi\ell}
\int_a^b\frac{\mathrm dr}r.
\]
故
\[
\boxed{R=\frac{\rho}{2\pi\ell}\ln\frac ba}.
\]
\excheck $b\to a$ 时电阻趋零；
不能把径向导电面积固定为内筒面积。

\electrocurrentproblem{155}{114}
半径为 $a$ 的金属球与半径为
$b>a$ 的同心金属薄球壳，
两者之间填满电阻率为 $\rho$
的均匀弱导电介质。
求径向电阻及 $b\to\infty$ 极限。
\exsolution
球对称径向电流的面积为
$A(r)=4\pi r^2$。
厚度 $\mathrm dr$ 的导电薄层电阻
$\mathrm dR=\rho\,\mathrm dr/(4\pi r^2)$。
积分：
\[
R=\frac{\rho}{4\pi}\int_a^b\frac{\mathrm dr}{r^2}
=\boxed{\frac{\rho}{4\pi}
\left(\frac1a-\frac1b\right)
=\frac{\rho(b-a)}{4\pi ab}}.
\]
极限
\[
\boxed{R_{b\to\infty}=\frac{\rho}{4\pi a}}.
\]
\excheck 无限介质中的球电极扩散电阻
仍为有限值，因远处径向导电面积按 $r^2$ 增长。

\electrocurrentproblem{156}{114}
半径为 $a<b$ 的同心球形电容器，
间隙填充均匀弱导电介质，
其电容为已知 $C$，
未接电源时电压在时间 $\Delta t$
内降低 $\eta>1$ 倍。
求介质电阻率 $\rho$。
\exsolution
球间电阻由第155题
$R=\rho(b-a)/(4\pi ab)$。
泄漏电流
$I=V/R$，
电荷满足
$\dot Q=-Q/(RC)$，
所以
$V(t)=V_0e^{-t/(RC)}$。
要求
$V_0/V(\Delta t)=\eta$，
即
$\Delta t/(RC)=\ln\eta$，
故
\[
R=\frac{\Delta t}{C\ln\eta}.
\]
联立球形电阻：
\[
\boxed{\rho=
\frac{4\pi ab\,\Delta t}
{(b-a)C\ln\eta}}.
\]
\excheck 原书简答的 PDF 提取可能
将 $\ln\eta$ 移入分子，
但指数衰减的时间常数
$\tau=\Delta t/\ln\eta$ 无歧义。

\electrocurrentproblem{157}{114}
相距远大于半径的两金属小球，
半径均为 $a$，
浸于电阻率 $\rho$ 的均匀弱导电介质中。
求两球之间的电阻。
\exsolution
两球之间电流通过周围导电介质传输。
对各球近场可视为单独球电极向远处导电，
一颗球的扩散电阻由上一题的极限为
$R_{\rm sph}=\rho/(4\pi a)$。
当两球相距 $\gg a$ 时，
两侧扩散电阻近似串联，
所以
\[
\boxed{R=2R_{\rm sph}=
\frac{\rho}{2\pi a}}.
\]
\excheck 两球若接近，
球面电流分布会显著畸变，
不能继续使用简单的两倍扩散电阻。

\electrocurrentproblem{158}{114}
半径为 $a$ 的金属小球距无限理想导体平面
$\ell\gg a$，周围填满电阻率为 $\rho$
的均匀弱导电介质。
若球与平面间电压为 $V$，
求 (a) 平面上电流密度分布；
(b) 两电极间总电阻。
\exsolution
稳态导电电位满足 Laplace 方程，
与同几何下的静电镜像法同构。
对小球近似一个点电流源，
其镜像为平面另一侧相反号点电流源，
两者在边界保持零电位。
因 $\ell\gg a$，总电流近似
$I\simeq4\pi aV/\rho$。
平面上距小球垂足水平距离为 $s$ 处，
两点电流源在边界的法向电流密度相加，
得到
\[
j_n(s)=\frac{I\ell}
{2\pi(s^2+\ell^2)^{3/2}}.
\]
因此
\[
\boxed{j_n(s)=
\frac{2a\ell V}
{\rho(s^2+\ell^2)^{3/2}}},
\qquad
\boxed{R_{\rm eff}\simeq\frac{\rho}{4\pi a}}.
\]
\excheck $2\pi\int_0^\infty s\,j_n(s)\mathrm ds=I$，
因此电流密度正确归一化；
原书使用的 $r$ 是从小球到平面点的距离
$r=\sqrt{s^2+\ell^2}$。

\electrocurrentproblem{159}{114}
两根无限长平行金属细线，
半径均为 $a$、轴间距为 $\ell\gg a$，
浸于电阻率 $\rho$ 的弱导电介质。
两线电势差为 $V$。
求 (a) 距两轴均为 $r$ 的点处电流密度模；
(b) 两线之间的单位长度导电电阻。
\exsolution
与二维线电荷电势同构，
两线沿轴单位长度电流记为 $I'$，
电势差在 $\ell\gg a$ 时为
\[
V\simeq\frac{\rho I'}{\pi}\ln\frac{\ell}{a}.
\]
因此每单位长度的\emph{电阻参数}
（其单位为 $\Omega\cdot\mathrm m$）为
\[
\boxed{\mathcal R'=V/I'
=\frac{\rho}{\pi}\ln\frac{\ell}{a}}.
\]
在距两线均为 $r$ 的中垂面上，
两个线电流源的场相加，
方向沿两轴连线，模为
$j=I'\ell/(2\pi r^2)$。
代入 $I'$：
\[
\boxed{j(r)=
\frac{\ell V}{2\rho r^2\ln(\ell/a)}}.
\]
\excheck “每单位长度的电阻”应理解为
$R\,L$（$\Omega\cdot$m），因为长度为
$L$ 的实际两线电阻为 $\mathcal R'/L$，
不是简单的 $\Omega/\mathrm m$。

\electrocurrentproblem{160}{114}
平板电容器板间填满电阻率
$\rho=100\,\mathrm{G\Omega\,m}$ 的玻璃，
电容 $C=4.0\,\mathrm{nF}$，
施加电压 $V=2.0\,\mathrm{kV}$。
求玻璃漏电流。
\exsolution
平板间距 $d$、面积 $S$ 时
$R=\rho d/S$，
$C=\varepsilon\varepsilon_0S/d$，
所以
\[
\boxed{RC=\rho\varepsilon\varepsilon_0}.
\]
漏电流
$I=V/R=VC/(\rho\varepsilon\varepsilon_0)$，
\[
\boxed{I=\frac{VC}{\rho\varepsilon\varepsilon_0}}.
\]
采用玻璃常用相对介电常数
$\varepsilon\simeq6$，
\[
\boxed{I\simeq
\frac{(2.0\times10^3)(4.0\times10^{-9})}
{(1.0\times10^{11})(6)(8.854\times10^{-12})}
\simeq1.5\,\mathrm{\mu A}}.
\]
\excheck $\varepsilon$ 未直接写在题目数据中，
数值需采用玻璃常用近似；
公式对其它玻璃类型仍适用。
